\documentclass[11pt,a4paper]{article}

\usepackage[margin=24mm]{geometry}
\usepackage{amsmath,amssymb,amsthm}
\usepackage{fontspec}
\usepackage{xeCJK}
\usepackage{microtype}
\usepackage[hidelinks]{hyperref}

\setmainfont{Latin Modern Roman}
\setCJKmainfont[AutoFakeSlant=0.2]{Noto Serif CJK SC}
\setCJKmonofont{Noto Sans Mono CJK SC}
\pagestyle{empty}
\setlength{\parindent}{0pt}
\setlength{\parskip}{0.7em}
\setlength{\jot}{7pt}

\newtheorem{lemma}{Lemma}
\newtheorem{proposition}{Proposition}
\newenvironment{paperproof}
  {\par\noindent\textit{Proof.}\par\smallskip}
  {\par}

\hypersetup{
  pdftitle={Fixed Points of a Random Permutation -- Reproduced English},
  pdfsubject={Unofficial reproduction for entertainment and educational purposes}
}

\begin{document}

\newcommand{\EditionName}{Reproduced English edition}
\newcommand{\EditionSummary}{%
  This edition preserves the text and calculation visible in the source
  video as closely as the available frames permit.}
\hypersetup{pageanchor=false}
\begin{titlepage}
\thispagestyle{empty}
\centering

\vspace*{1.2cm}

{\Huge\bfseries Fixed Points of a Random Permutation\par}
\vspace{0.45cm}
{\Large Reproduction and Revised Editions\par}

\vspace{1.1cm}

\begin{minipage}{0.86\textwidth}
\raggedright

\textbf{Source.}
This document is an independent LaTeX reproduction of, and editorial
adaptation based on, the mathematical paper shown in the Bilibili video
\emph{【日常研究】摆烂小能手掀起的数学问题} (video ID
\texttt{BV1Fh411N71S}):
\begin{center}
  \href{https://www.bilibili.com/video/BV1Fh411N71S}
       {\nolinkurl{https://www.bilibili.com/video/BV1Fh411N71S}}
\end{center}

\textbf{Repository.}
The \LaTeX\ sources, compiled PDFs, and edition notes are published at:
\begin{center}
  \href{https://github.com/Ritsuka314/personal-website}
       {\nolinkurl{https://github.com/Ritsuka314/personal-website}}
\end{center}

\textbf{PDF editions.}
\begin{center}
  \small
  \href{https://opuscula.ritsuka.moe/euler-derangement-cloze/paper-en.pdf}
       {\texttt{paper-en.pdf}}
  \quad
  \href{https://opuscula.ritsuka.moe/euler-derangement-cloze/paper-en-revised.pdf}
       {\texttt{paper-en-revised.pdf}}
  \quad
  \href{https://opuscula.ritsuka.moe/euler-derangement-cloze/paper-zh.pdf}
       {\texttt{paper-zh.pdf}}
\end{center}

\textbf{Purpose.}
The reconstruction was prepared solely for entertainment and educational
purposes. It is unofficial, is not presented as an author-approved
transcript, and makes no claim of affiliation with the video creator or
platform.

\medskip
\textbf{The three editions.}
\begin{enumerate}
  \item \textbf{Reproduced English edition.}
    A close typesetting of the four-page English note visible in the video.
    It intentionally retains the source's nonstandard wording, omissions,
    notation, and induction argument.

  \item \textbf{Revised English edition.}
    A corrected, self-contained mathematical report. It defines the model
    precisely, proves the derangement formula by inclusion--exclusion, gives
    a unified probability distribution, and proves the expectation with
    indicator variables.

  \item \textbf{Revised Chinese edition.}
    A polished Chinese counterpart to the revised English report. It follows
    the same definitions, results, and proofs rather than translating the
    reproduced wording literally.
\end{enumerate}

\medskip
\noindent
\fbox{%
  \begin{minipage}{0.93\linewidth}
    \textbf{Edition in this file: \EditionName}\par
    \EditionSummary
  \end{minipage}%
}

\end{minipage}

\vfill
{\small Typeset with \LaTeX\ from the source video cited above.\par}

\end{titlepage}
\hypersetup{pageanchor=true}


\begin{lemma}
Consider a permutation of \(n\) elements:
\[
  f:\{1,2,\ldots,n\}\longrightarrow\{1,2,\ldots,n\},
\]
where \(f\) is bijective. Then the probability of getting the permutation such
that
\[
  f(k)\ne k\bigl(\forall k\in\{1,2,\ldots,n\}\bigr)
\]
is
\[
  \sum_{i=2}^{n}\frac{(-1)^i}{i!}.
\]
And the total number of such \(f\) is
\[
  n!\sum_{i=2}^{n}\frac{(-1)^i}{i!}.
\]
\end{lemma}

The proof of this lemma is omitted here. Assume \(X\) is a random variable
referring to how many elements in the permutation are corrected aligned
(i.e.\ \(X=\) the number of \(k\)'s such that \(f(k)=k\)). Our problem is to
calculate the probability distribution of \(X\) and its expectation.

\newpage

The trivial cases are
\[
  P(X=n)=\frac{1}{n!},
\]
and
\[
  P(X=n-1)=0.
\]

In general, for \(k\in\{0,1,2,\ldots,n-2\}\),
\begin{align*}
  P(X=k)
    &=P\bigl(k\text{ elements aligned with the other }(n-k)
      \text{ elements ALL misaligned}\bigr)\\
    &=\frac{1}{n!}\binom{n}{k}
      \left(\sum_{i=2}^{n-k}\frac{(-1)^i}{i!}\right)(n-k)!\\
    &=\frac{1}{k!}\sum_{i=2}^{n-k}\frac{(-1)^i}{i!}.
\end{align*}

where \(n!\) is the total number of cases considered, \(\binom{n}{k}\) is the
binomial coefficient \(C_n^k\), and
\[
  \left(\sum_{i=2}^{n-k}\frac{(-1)^i}{i!}\right)(n-k)!
\]
is the number of cases where all \((n-k)\) fixed elements are misaligned in
the alignment.

Therefore, the probability distribution of \(X\) is
\[
  P(X=k)=
  \begin{cases}
    \displaystyle\frac{1}{k!}\sum_{i=2}^{n-k}\frac{(-1)^i}{i!},
      & \text{if }k\in\{0,1,2,\ldots,n-2\},\\[1.1em]
    0, & \text{if }k=n-1,\\[0.25em]
    \displaystyle\frac{1}{n!}, & \text{if }k=n.
  \end{cases}
\]

\newpage

\begin{proposition}
The expectation of \(X\) is
\[
  E(X)=1.
\]
\end{proposition}

\begin{paperproof}
\begin{align*}
  E(X)
    &=\frac{n}{n!}+(n-1)\cdot 0
      +\sum_{k=0}^{n-2}\frac{k}{k!}
       \sum_{i=2}^{n-k}\frac{(-1)^i}{i!}\\
    &=\frac{1}{(n-1)!}
      +\sum_{k=1}^{n-2}\frac{1}{(k-1)!}
       \sum_{i=2}^{n-k}\frac{(-1)^i}{i!}\\
    &=\frac{1}{(n-1)!}
      +\sum_{k=1}^{n-2}\sum_{i=2}^{n-k}
       \frac{(-1)^i}{i!(k-1)!}.
\end{align*}

Use induction on \(n\). Suppose we already have
\[
  \frac{1}{(n-1)!}
  +\sum_{k=1}^{n-2}\sum_{i=2}^{n-k}
   \frac{(-1)^i}{i!(k-1)!}=1.
\]

\newpage

It follows that, for the \((n+1)\) case,
\begingroup
\small
\setlength{\jot}{6pt}
\begin{align*}
  &\frac{1}{n!}
   +\sum_{k=1}^{n-1}\sum_{i=2}^{n+1-k}
    \frac{(-1)^i}{i!(k-1)!}\\
  ={}&\frac{1}{n!}+\frac{1}{2\times(n-2)!}
   +\sum_{k=1}^{n-2}\sum_{i=2}^{n+1-k}
    \frac{(-1)^i}{i!(k-1)!}\\
  ={}&\frac{1}{n!}+\frac{1}{2\times(n-2)!}
   +\sum_{k=1}^{n-2}\sum_{i=2}^{n-k}
    \frac{(-1)^i}{i!(k-1)!}
   +\sum_{k=1}^{n-2}
    \frac{(-1)^{n+1-k}}{(n+1-k)!(k-1)!}\\
  ={}&\frac{1}{n!}+\frac{1}{2\times(n-2)!}
   +1-\frac{1}{(n-1)!}
   +\sum_{k=1}^{n-2}
    \frac{(-1)^{n+1-k}}{(n+1-k)!(k-1)!}\\
  ={}&\frac{n^2-3n+2}{2\times n!}+1
   +\sum_{k=1}^{n-2}
    \frac{(-1)^{n+1-k}}{(n+1-k)!(k-1)!}\\
  ={}&\frac{n^2-3n+2}{2\times n!}+1
   +\sum_{t=0}^{n-3}\frac{(-1)^{n-t}}{(n-t)!t!}\\
  ={}&\frac{n^2-3n+2}{2\times n!}+1
   +\frac{1}{n!}\sum_{t=0}^{n-3}(-1)^{n-t}\binom{n}{t}\\
  ={}&\frac{n^2-3n+2}{2\times n!}+1
   -\frac{1}{n!}(-1)^2\binom{n}{n-2}
   -\frac{1}{n!}(-1)^1\binom{n}{n-1}\\[-0.1em]
   &\quad
   -\frac{1}{n!}(-1)^0\binom{n}{n}
   +\frac{1}{n!}\sum_{t=0}^{n}(-1)^{n-t}\binom{n}{t}\\
  ={}&\frac{n^2-3n+2}{2\times n!}+1
   -\frac{1}{n!}\binom{n}{2}
   +\frac{1}{n!}\binom{n}{1}
   -\frac{1}{n!}
   +\frac{1}{n!}(-1+1)^n\\
  ={}&\frac{n^2-3n+2-n(n-1)+2n-2}{2\times n!}+1\\
  ={}&1.
\end{align*}
\endgroup

By induction, the result
\[
  E(X)=1
\]
holds for all \(n\in\mathbb{N}\).
\end{paperproof}

\end{document}
